% --- Para probar este código de forma independiente ---
\documentclass[tikz, border=10pt]{standalone}

% --- Paquetes necesarios ---
\usepackage[utf8]{inputenc}
\usepackage[T1]{fontenc}
\usepackage{lmodern} % Mejores fuentes
\usepackage{amsmath} % Para el nodo de texto
\usepackage{graphicx}
% --- Librerías TikZ ---
\usetikzlibrary{
    patterns.meta,  % Para el achurado (hatch) del pistón
    arrows.meta,    % Para las flechas (mejor que 'arrows')
    calc            % Para cálculos de coordenadas
}

\begin{document}
 

\begin{tikzpicture}[
    % --- Estilos Globales ---
    thick_wall/.style={line width=1.5pt},
    piston_pattern/.style={
        pattern={Lines[angle=45, distance=3pt, line width=0.8pt]}
    },
    big_arrow/.style={
        -Triangle [length=8mm, width=8mm] , % Flecha grande y gruesa
        line width=2.5pt
    },
    pressure_arrow/.style={
        -Stealth, % Flecha de presión estándar
        thick
    }
]

% ==========================================================
% 1. DEFINICIÓN DE DIMENSIONES (Paramétrico)
% ==========================================================
\def\W{6}      % Ancho interno del cilindro
\def\T{0.4}    % Grosor de las paredes del cilindro
\def\Hbase{6}  % Altura inicial del gas (de A-B a C-D)
\def\Hpist{1}  % Grosor/Altura del pistón
\def\Hrod{5}   % Altura del vástago (mango)
\def\Wrod{1}   % Ancho del vástago (mango)
\def\dL{4}     % Desplazamiento (distancia dL)
\def\Htop{7}   % Altura total de la pared interna (debe ser > Hbase+dL+Hrod)

% ==========================================================
% 2. DIBUJO DEL CILINDRO
% ==========================================================
% --- Coordenadas base internas ---
\coordinate (A) at (-\W/2, 0);
\coordinate (B) at (\W/2, 0);

% --- Paredes (dibujadas como cajas) ---
% Relleno blanco para tapar cualquier cosa detrás (si se usara \includegraphics)
\fill[white] (-\W/2-\T, -\T) rectangle (\W/2+\T, \Hbase+\Htop);

% --- Líneas de las paredes ---
\draw[thick_wall] (A) -- (-\W/2, \Hbase+\Htop); % Pared interna izq
\draw[thick_wall] (B) -- (\W/2, \Hbase+\Htop);  % Pared interna der
\draw[thick_wall] (A) -- (B);                     % Base interna A-B

\draw[thick_wall] (-\W/2-\T, -\T) -- (-\W/2-\T, \Hbase+\Htop); % Pared externa izq
\draw[thick_wall] (\W/2+\T, -\T) -- (\W/2+\T, \Hbase+\Htop);  % Pared externa der
\draw[thick_wall] (-\W/2-\T, -\T) -- (\W/2+\T, -\T);          % Base externa

% ==========================================================
% 3. SISTEMA DE GAS (Posición inicial)
% ==========================================================
% --- Coordenadas de la parte superior del gas ---
\coordinate (C) at (-\W/2, \Hbase);
\coordinate (D) at (\W/2, \Hbase);

% --- Caja punteada "Límite del Sistema" ---
\draw[dashed, thick] (A) -- (B) -- (D) -- (C) -- cycle;
\node at (0, \Hbase/2) {\Huge GAS};

% --- Etiqueta "Límite del Sistema" ---
\draw[-Stealth, thick] (\W/2 + 2.5, \Hbase*0.7) -- ($(B) + (0.5, 0.5)$);
\node[right, align=left] at (\W/2 + 2.5, \Hbase*0.7) {LIMITE DEL \\ SISTEMA};

% ==========================================================
% 4. PISTÓN Y VÁSTAGO (Posición inicial)
% ==========================================================
% --- Pistón (achurado) ---
\fill[piston_pattern] (C) rectangle (\W/2, \Hbase + \Hpist);
\draw[thick_wall] (C) rectangle (\W/2, \Hbase + \Hpist);

% --- Vástago (mango) ---
\coordinate (Rod_BL) at (-\Wrod/2, \Hbase + \Hpist); % Base Izq
\coordinate (Rod_TR) at (\Wrod/2, \Hbase + \Hpist + \Hrod); % Arriba Der
\draw[thick_wall, fill=white] (Rod_BL) rectangle (Rod_TR);
\node[right, xshift=3mm] at (Rod_TR.north west) {$p$};

% --- Flechas de Presión (sobre el pistón, no sobre el vástago) ---
\foreach \x in {-\W/2+0.7, -\W/2+1.8, \W/2-0.7, \W/2-1.8} {
    \draw[pressure_arrow] (\x, \Hbase + \Hpist + 0.7) -- (\x, \Hbase + \Hpist);
}

% ==========================================================
% 5. POSICIÓN FINAL (Punteada) Y COTA "dL"
% ==========================================================
% --- Coordenadas de la posición final ---
\coordinate (E) at (-\W/2, \Hbase + \dL);
\coordinate (F) at (\W/2, \Hbase + \dL);

% --- Dibujo punteado del pistón y vástago ---
\draw[dashed, thick] (E) rectangle (\W/2, \Hbase + \dL + \Hpist);
\draw[dashed, thick] ($(Rod_BL) + (0, \dL)$) rectangle ($(Rod_TR) + (0, \dL)$);

% --- Cota dL ---
\draw[<->, thick] (\W/2 + \T + 0.7, \Hbase) -- (\W/2 + \T + 0.7, \Hbase + \dL) 
    node[midway, right, xshift=2mm] {$dL$};

% ==========================================================
% 6. FLECHAS DE ENERGÍA Y ETIQUETAS A-F
% ==========================================================
% --- Energía Calorífica (Abajo) ---
\draw[big_arrow] (0, -2.5) -- (0, -\T) 
    node[midway, left, xshift=-5mm, align=center] {ENERGIA \\ CALORIFICA};

% --- Energía Mecánica (Arriba) ---
% Se origina en la posición final del vástago
\coordinate (Rod_Top_Final) at (0, \Hbase + \Hpist + \Hrod + \dL);
\draw[big_arrow] (Rod_Top_Final) -- ($(Rod_Top_Final) + (0, 4.5)$) 
    node[midway, right, xshift=5mm, align=center] {ENERGIA \\ MECANICA};

% --- Etiquetas A, B, C, D, E, F ---
%\node[left=2mm of A] {A};
%\node[right=2mm of B] {B};
%\node[left=2mm of C] {C};
%\node[right=2mm of D] {D};
%\node[left=2mm of E] {E};
%\node[right=2mm of F] {F};

\end{tikzpicture}

\end{document}